\section{The rational function field}

Let $k$ be an arbitrary field, put $R=k[T]$ and $F=k(T)$, and let
$\Sigma$ be the set of nonzero prime ideals of $R$. The parameter set
will be $\Omega=\{\infty\}\sqcup\Sigma$, with its discrete sigma algebra.
No countability assumption on $k$ or $\Omega$ is made.

Write $\Gamma=\{0\}\cup\{\tau^m:m\in\mathbb Z\}$ for the ordered
multiplicative group of integers with zero adjoined, where
$\tau^m\tau^n=\tau^{m+n}$ and $\tau^m\leq\tau^n$ if and only if
$m\leq n$. We use the standard valuations of rational functions and the
factorization of ideals in the principal ideal domain $k[T]$, in the
forms recorded in \cite{mathlib-rational-function-valuations}.
Explicitly, each $v\in\Sigma$ gives a valuation $V_v:F\to\Gamma$
whose restriction $V_v^R$ to $R$ satisfies
\[
 V_v^R(r)\leq1,\qquad
 V_v^R(r)=1\ \Longleftrightarrow\ r\notin v,\qquad
 V_v^R(r)<1\ \Longleftrightarrow\ v\mid(r).
\]
If $v=(q)$ with $q\ne0$, its normalization is $V_v^R(q)=\tau^{-1}$.
The valuation at infinity satisfies
$V_\infty(p/q)=\tau^{\deg p-\deg q}$ for nonzero polynomials $p,q$.
These valuations are multiplicative, vanish exactly at zero, and satisfy
the ultrametric inequality. Every nonzero ideal of $R$ has only finitely
many prime-ideal divisors. These are the valuation and factorization
inputs; the product formula below is proved from them.

\begin{definition}[Exponential base]\label{def:ff-exp-base}
\lean{ArakelovGeometry.AdelicStructure.expBase,ArakelovGeometry.AdelicStructure.one_lt_expBase,ArakelovGeometry.AdelicStructure.coe_expBase}
Let $e=\exp(1)$, regarded as a nonnegative real number. Then its underlying
real value is $\exp(1)$ and $e>1$.
\end{definition}
\begin{proof}
The exponential is positive, and its strict monotonicity gives
$\exp(1)>\exp(0)=1$. Thus regarding it as a nonnegative real number
does not change its real value.
\end{proof}

\begin{lemma}[Logarithm of the exponential realization]
\label{lem:ff-log-exponential}
\uses{def:ff-exp-base}
\lean{ArakelovGeometry.AdelicStructure.real_log_toNNReal_expBase}
Define $\lambda:\Gamma\to\mathbb R_{\geq0}$ by $\lambda(0)=0$ and
$\lambda(\tau^m)=e^m$. For $\gamma=\tau^m\ne0$, one has
$\log\lambda(\gamma)=m$.
\end{lemma}
\begin{proof}
Since $e>0$, the logarithm of an integer power is its exponent times
$\log e$. The latter logarithm is $1$, so $\log(e^m)=m$.
\end{proof}

\begin{definition}[Degree of a finite point]\label{def:ff-point-degree}
\lean{ArakelovGeometry.AdelicStructure.pointDegree}
For $v\in\Sigma$, choose a generator $g_v$ of the principal ideal $v$
and define $d_v=\deg g_v\in\mathbb N$.
\end{definition}

\begin{lemma}[Degree is independent of a generator]
\label{lem:ff-generator-degree}
\uses{def:ff-point-degree}
\lean{ArakelovGeometry.AdelicStructure.pointDegree_eq_of_span}
If $v=(q)$, then $d_v=\deg q$.
\end{lemma}
\begin{proof}
The generators $g_v$ and $q$ are nonzero because $v\ne(0)$.
Equality of their principal ideals implies $g_v=u q$ for a unit $u$ of
$k[T]$. Every such unit is a nonzero constant, of degree zero.
The degree formula for a product of nonzero polynomials therefore gives
$\deg g_v=\deg u+\deg q=\deg q$.
\end{proof}

\begin{definition}[The point of a prime polynomial]\label{def:ff-prime-point}
\lean{ArakelovGeometry.AdelicStructure.spanPrime,ArakelovGeometry.AdelicStructure.spanPrime_asIdeal}
For a prime polynomial $q\in k[T]$, let $v_q\in\Sigma$ be the point
whose ideal is $(q)$.
\end{definition}
\begin{proof}
A prime polynomial is nonzero and not a unit. Its principal ideal is
therefore a nonzero proper prime ideal, as required.
\end{proof}

\begin{lemma}[Degree of a prime-polynomial point]\label{lem:ff-prime-degree}
\uses{def:ff-prime-point,lem:ff-generator-degree}
\lean{ArakelovGeometry.AdelicStructure.pointDegree_spanPrime}
For a prime polynomial $q$, one has $d_{v_q}=\deg q$.
\end{lemma}
\begin{proof}
The ideal of $v_q$ is $(q)$, so the generator-degree formula applies.
\end{proof}

\begin{lemma}[Logarithm at a finite point]\label{lem:ff-finite-log}
\uses{lem:ff-log-exponential}
\lean{ArakelovGeometry.AdelicStructure.real_log_adicAbv}
Let $|f|_v=\lambda(V_v(f))$ be the $v$-adic absolute value with base
$e$. If $f\ne0$ and $V_v(f)=\tau^m$, then $\log|f|_v=m$.
\end{lemma}
\begin{proof}
The valuation of a nonzero element is nonzero. Apply
Lemma~\ref{lem:ff-log-exponential} to $V_v(f)$.
\end{proof}

\begin{lemma}[Valuations of a prime polynomial]\label{lem:ff-prime-valuations}
\uses{def:ff-prime-point}
\lean{ArakelovGeometry.AdelicStructure.valuation_spanPrime_self,ArakelovGeometry.AdelicStructure.valuation_algebraMap_eq_one_of_ne}
For a prime polynomial $q$, regarded as an element of $F$,
\[
 V_{v_q}(q)=\tau^{-1},\qquad
 V_v(q)=1\quad(v\ne v_q).
\]
\end{lemma}
\begin{proof}
The first assertion is the normalization of the valuation at its
principal prime ideal. For the second, suppose $q\in v$. Then
$(q)\subseteq v$. A nonzero prime ideal of $k[T]$ is maximal, so this
inclusion of proper prime ideals forces $v=(q)$, contrary to $v\ne v_q$.
Thus $q\notin v$, and the valuation criterion gives $V_v^R(q)=1$.
The valuation on $F$ restricts to $V_v^R$, proving the assertion.
\end{proof}

\begin{lemma}[Valuations of nonzero constants]\label{lem:ff-constant-valuation}
\lean{ArakelovGeometry.AdelicStructure.valuation_algebraMap_C}
For $c\in k^\times$ and $v\in\Sigma$, one has $V_v(c)=1$.
\end{lemma}
\begin{proof}
The constant polynomial $c$ is a unit of $k[T]$ and cannot belong to
the proper ideal $v$. The criterion $V_v^R(c)=1$ if and only if
$c\notin v$, together with restriction of the valuation to $R$, proves
the assertion.
\end{proof}

\begin{definition}[Absolute value at infinity]\label{def:ff-infinity-absolute-value}
\uses{def:ff-exp-base,lem:ff-log-exponential}
\lean{ArakelovGeometry.AdelicStructure.inftyAbvDef,ArakelovGeometry.AdelicStructure.isNonarchimedean_inftyAbvDef,ArakelovGeometry.AdelicStructure.inftyAbv}
Put $|f|_\infty=\lambda(V_\infty(f))$. This is a nonarchimedean real
absolute value on $F$. It takes the value zero at zero and, for nonzero
$p,q\in k[T]$, satisfies $|p/q|_\infty=e^{\deg p-\deg q}$.
\end{definition}
\begin{proof}
The map $\lambda$ is multiplicative and strictly increasing: these
properties follow from the exponent laws and $e>1$. It takes zero only
at zero. Composition with $V_\infty$ therefore preserves
multiplicativity, definiteness, and the ultrametric inequality. The
ultrametric inequality implies the ordinary triangle inequality because
the values are nonnegative. The displayed formula follows from the
formula for $V_\infty(p/q)$.
\end{proof}

\begin{lemma}[Logarithmic degree at infinity]\label{lem:ff-infinity-polynomial-log}
\uses{def:ff-infinity-absolute-value,lem:ff-log-exponential}
\lean{ArakelovGeometry.AdelicStructure.real_log_inftyAbv_algebraMap}
For every nonzero polynomial $p\in k[T]$, one has
$\log|p|_\infty=\deg p$.
\end{lemma}
\begin{proof}
The valuation at infinity of $p$ is $\tau^{\deg p}$. Its logarithm
after applying $\lambda$ is $\deg p$ by
Lemma~\ref{lem:ff-log-exponential}.
\end{proof}

\begin{definition}[Absolute values and weights on the parameter set]
\label{def:ff-point-data}
\uses{def:ff-point-degree,lem:ff-finite-log,def:ff-infinity-absolute-value}
\lean{ArakelovGeometry.AdelicStructure.pointAbv,ArakelovGeometry.AdelicStructure.pointWeight,ArakelovGeometry.AdelicStructure.pointAbv_none,ArakelovGeometry.AdelicStructure.pointAbv_some,ArakelovGeometry.AdelicStructure.pointWeight_none,ArakelovGeometry.AdelicStructure.pointWeight_some}
On $\Omega=\{\infty\}\sqcup\Sigma$, use the absolute value
$|\cdot|_\infty$ at $\infty$ and $|\cdot|_v$ at a finite point $v$.
Set $n_\infty=1$ and $n_v=d_v$, regarded as elements of $[0,+\infty]$.
For $f\in F$, write
\[
 \ell_f(\omega)=\log|f|_\omega,\qquad
 h_f(\omega)=n_\omega\ell_f(\omega),
\]
where the weight in the second formula is regarded as a real number
and $\log0=0$.
\end{definition}

\begin{lemma}[Finiteness of the weights]\label{lem:ff-finite-weights}
\uses{def:ff-point-data}
\lean{ArakelovGeometry.AdelicStructure.pointWeight_ne_top}
For every $\omega\in\Omega$, one has $n_\omega<+\infty$.
\end{lemma}
\begin{proof}
At infinity the weight is $1$; at a finite point it is the natural
number $d_v$. Both are finite.
\end{proof}

\begin{lemma}[Clearing the canonical denominator]\label{lem:ff-clearing-denominators}
\lean{ArakelovGeometry.AdelicStructure.algebraMap_num_eq_mul_denom}
Write $a_f,b_f\in k[T]$ for the canonical numerator and denominator of
$f\in F$, with $b_f\ne0$. Then $a_f=f b_f$ in $F$.
\end{lemma}
\begin{proof}
The canonical fraction identity is $f=a_f/b_f$.
The polynomial embedding into $F$ is injective, so $b_f$ remains
nonzero in $F$. Multiplication by it gives the assertion.
\end{proof}

\begin{lemma}[The local unit criterion]\label{lem:ff-local-unit-criterion}
\lean{ArakelovGeometry.AdelicStructure.intValuation_eq_one_of_not_dvd}
If $r\in k[T]$ and $v\nmid(r)$, then $V_v^R(r)=1$.
\end{lemma}
\begin{proof}
One always has $V_v^R(r)\leq1$. If its value were strictly less than
$1$, the valuation-divisibility criterion would give $v\mid(r)$.
The hypothesis thus forces equality.
\end{proof}

\begin{lemma}[Unit valuation of a quotient]\label{lem:ff-quotient-unit-valuation}
\uses{lem:ff-clearing-denominators}
\lean{ArakelovGeometry.AdelicStructure.valuation_eq_one_of_intValuation_num_denom}
If $V_v^R(a_f)=V_v^R(b_f)=1$, then $V_v(f)=1$.
\end{lemma}
\begin{proof}
Apply $V_v$ to $a_f=f b_f$. Multiplicativity and restriction to $R$
give $1=V_v(f)\cdot1$, hence the result.
\end{proof}

\begin{lemma}[Support of logarithmic absolute values]\label{lem:ff-support-bound}
\uses{def:ff-point-data,lem:ff-finite-log,lem:ff-local-unit-criterion,lem:ff-quotient-unit-valuation}
\lean{ArakelovGeometry.AdelicStructure.support_log_pointAbv_subset}
For every $f\in F$,
\[
 \operatorname{supp}(\ell_f)\subseteq
 \{\infty\}\cup\{v\in\Sigma:v\mid(a_f)\text{ or }v\mid(b_f)\}.
\]
\end{lemma}
\begin{proof}
For $f=0$, every absolute value vanishes and the convention $\log0=0$
makes the support empty. Suppose $f\ne0$, and let a finite point $v$
divide neither ideal. The local unit criterion gives
$V_v^R(a_f)=V_v^R(b_f)=1$. Hence $V_v(f)=1$ by
Lemma~\ref{lem:ff-quotient-unit-valuation}, and the logarithmic formula
at a finite point gives $\ell_f(v)=0$. This proves the inclusion.
\end{proof}

\begin{lemma}[Finite logarithmic support]\label{lem:ff-finite-log-support}
\uses{lem:ff-support-bound}
\lean{ArakelovGeometry.AdelicStructure.finite_support_log}
For every $f\in F^\times$, the function $\ell_f$ has finite support.
\end{lemma}
\begin{proof}
The numerator $a_f$ and denominator $b_f$ are nonzero. Their principal
ideals are therefore nonzero and have only finitely many prime-ideal
divisors. The support bound is a subset of the union of these two finite
sets and the singleton $\{\infty\}$.
\end{proof}

\begin{lemma}[Support of the weighted logarithm]\label{lem:ff-weighted-log-support}
\uses{def:ff-point-data,lem:ff-finite-log-support}
\lean{ArakelovGeometry.AdelicStructure.support_weighted_log_subset,ArakelovGeometry.AdelicStructure.finite_support_weighted_log}
For every $f\in F$, one has
$\operatorname{supp}(h_f)\subseteq\operatorname{supp}(\ell_f)$.
In particular, $h_f$ has finite support if $f\ne0$.
\end{lemma}
\begin{proof}
If $\ell_f(\omega)=0$, multiplication by the real weight gives
$h_f(\omega)=0$. The asserted support inclusion follows; the finiteness
claim follows from Lemma~\ref{lem:ff-finite-log-support}.
\end{proof}

For $f\in F^\times$, define
$B(f)=\sum_{\omega\in\Omega}^{\mathrm{fin}}h_f(\omega)$,
where the finite sum may be taken over any finite set containing the
support of $h_f$.

\begin{lemma}[Additivity of the weighted logarithmic sum]
\label{lem:ff-product-log-additive}
\uses{def:ff-point-data,lem:ff-weighted-log-support}
\lean{ArakelovGeometry.AdelicStructure.finsum_weighted_log_mul}
For $f,g\in F^\times$, one has $B(fg)=B(f)+B(g)$.
\end{lemma}
\begin{proof}
Every absolute value is positive on $f$ and $g$. Multiplicativity and
the logarithm law therefore give
$\ell_{fg}(\omega)=\ell_f(\omega)+\ell_g(\omega)$.
Multiplying by the weight gives $h_{fg}=h_f+h_g$. All three functions
have finite support, so summing over the union of these supports proves
the equality.
\end{proof}

\begin{lemma}[Product formula for constants]\label{lem:ff-constant-product-formula}
\uses{def:ff-point-data,lem:ff-constant-valuation,lem:ff-finite-log,lem:ff-infinity-polynomial-log}
\lean{ArakelovGeometry.AdelicStructure.finsum_weighted_log_C}
For $c\in k^\times$, one has $B(c)=0$.
\end{lemma}
\begin{proof}
At every finite point, $V_v(c)=1$, so $\log|c|_v=0$. At infinity the
logarithm is the degree of the nonzero constant polynomial $c$, also
zero. Thus every summand in $B(c)$ is zero.
\end{proof}

\begin{lemma}[Product formula for prime polynomials]
\label{lem:ff-prime-product-formula}
\uses{def:ff-point-data,lem:ff-prime-degree,lem:ff-prime-valuations,lem:ff-finite-log,lem:ff-infinity-polynomial-log}
\lean{ArakelovGeometry.AdelicStructure.finsum_weighted_log_prime}
For a prime polynomial $q\in k[T]$, one has $B(q)=0$.
\end{lemma}
\begin{proof}
At any finite point other than $v_q$, the valuation is $1$, so its
logarithm is zero. At $v_q$ it is $\tau^{-1}$ and the logarithm is
$-1$; its weight is $d_{v_q}=\deg q$. At infinity the logarithm is
$\deg q$ and the weight is $1$. Consequently the sum reduces to
$-(\deg q)+\deg q=0$.
\end{proof}

\begin{lemma}[Product formula for polynomials]\label{lem:ff-polynomial-product-formula}
\uses{lem:ff-constant-product-formula,lem:ff-prime-product-formula,lem:ff-product-log-additive}
\lean{ArakelovGeometry.AdelicStructure.finsum_weighted_log_polynomial}
For every nonzero polynomial $p\in k[T]$, one has $B(p)=0$.
\end{lemma}
\begin{proof}
Unique factorization writes $p=cq_1\cdots q_r$, where $c\in k^\times$
and the $q_i$ are prime polynomials, listed with multiplicity.
Repeated additivity gives
$B(p)=B(c)+\sum_{i=1}^r B(q_i)$. The constant and prime-polynomial
formulas show that each term is zero, including the case $r=0$.
\end{proof}

\begin{theorem}[Product formula for rational functions]\label{thm:ff-product-formula}
\uses{lem:ff-clearing-denominators,lem:ff-polynomial-product-formula,lem:ff-product-log-additive}
\lean{ArakelovGeometry.AdelicStructure.finsum_weighted_log_eq_zero}
For every $f\in F^\times$,
\[
 \sum_{\omega\in\Omega}^{\mathrm{fin}}n_\omega\log|f|_\omega=0.
\]
\end{theorem}
\begin{proof}
The numerator and denominator of $f$ are nonzero, and
$a_f=f b_f$. Applying additivity and the polynomial formula gives
$0=B(a_f)=B(f)+B(b_f)=B(f)$.
\end{proof}

\begin{definition}[The rational-function adelic curve]\label{def:ff-adelic-structure}
\uses{def:adelic-structure,def:singleton-measure,lem:singleton-measure-values,def:ff-point-data,lem:ff-finite-weights,lem:ff-finite-log-support,lem:finite-support-integrability}
\lean{ArakelovGeometry.AdelicStructure.functionFieldAdelicStructure,ArakelovGeometry.AdelicCurve.functionFieldAdelicCurve}
Equip the discrete measurable space $\Omega$ with the measure
$\nu=\sum_{\omega\in\Omega}n_\omega\delta_\omega$ and the absolute
values $|\cdot|_\omega$. These data define an adelic structure on
$F=k(T)$, and hence an adelic curve.
\end{definition}
\begin{proof}
Every function on $\Omega$ is measurable because its sigma algebra is
discrete. The measure of a singleton $\{\omega\}$ is $n_\omega$, which
is finite. For $f\ne0$, its logarithmic absolute-value function
$\ell_f$ has finite support by Lemma~\ref{lem:ff-finite-log-support}.
Lemma~\ref{lem:finite-support-integrability} therefore makes $\ell_f$
integrable. These are exactly the conditions for an adelic structure.
\end{proof}

\begin{theorem}[The integral product formula and properness]
\label{thm:ff-integral-product-formula}
\uses{def:ff-adelic-structure,lem:singleton-measure-values,lem:ff-finite-weights,lem:ff-finite-log-support,lem:ff-weighted-log-support,lem:finite-support-integral,thm:ff-product-formula}
\lean{ArakelovGeometry.AdelicCurve.functionField_product_formula}
For $f\in F^\times$, one has
$\int_\Omega\log|f|_\omega\,d\nu(\omega)=0$.
In particular, the rational-function adelic curve is proper.
\end{theorem}
\begin{proof}
Choose a finite set $A\subseteq\Omega$ containing the support of
$\ell_f$. Every singleton in $A$ has finite measure, so
Lemma~\ref{lem:finite-support-integral} gives
\[
 \int_\Omega\ell_f\,d\nu
 =\sum_{\omega\in A}\nu(\{\omega\})\ell_f(\omega)
 =\sum_{\omega\in A}n_\omega\ell_f(\omega).
\]
The support of the weighted logarithm is contained in that of
$\ell_f$, so the last sum is $B(f)$. This is zero by
Theorem~\ref{thm:ff-product-formula}. The defining product formula for
properness holds for every nonzero $f$, proving the last assertion.
\end{proof}
